% 来源及取材范围见分题文件；此为整理者转录的正文，不是下载原件。
\paragraph{Definitions from \S5.1 and Definition 5.30 (collected notation).}
\noindent\textit{以下背景与记号据所列原文章节摘编；不是逐字引文，也不是新增证明。}\par
A graphon is a symmetric measurable function $W:[0,1]^2\to[0,1]$. The cut norm is
\[\|W\|_\square=\sup_{S,T\subseteq[0,1]}\left|\int_{S\times T}W\right|,\]
where $S,T$ are measurable. Write $W^\varphi(x,y)=W(\varphi(x),\varphi(y))$ for a measure-preserving bijection $\varphi$ of $[0,1]$ and
\[\delta_\square(U,W)=\inf_\varphi\|U-W^\varphi\|_\square.\]
Let $\widetilde{\mathcal W}_0$ be the space of graphons modulo $U\sim W$ if $\delta_\square(U,W)=0$. For a finite measurable partition $\mathcal P=\{S_1,\ldots,S_k\}$, the stepping operator is
\[W_{\mathcal P}(x,y)=\frac1{\lambda(S_i)\lambda(S_j)}\int_{S_i\times S_j}W\quad ((x,y)\in S_i\times S_j).\]
Sets of measure zero are ignored. This is projection in $L^2([0,1]^2)$ onto the functions constant on each step, or conditional expectation with respect to the generated $\sigma$-algebra.
\begin{lemma}[5.34, $L^2$ energy increment]
Let $W$ be a graphon and $\mathcal P$ a partition of $[0,1]$, satisfying $\|W-W_{\mathcal P}\|_\square>\epsilon$. There exists a refinement $\mathcal P'$ of $\mathcal P$ dividing each part into no more than $4$ parts, such that
\[\|W_{\mathcal P'}\|_2^2>\|W_{\mathcal P}\|_2^2+\epsilon^2.\]
\end{lemma}
\begin{proof}
Because $\|W-W_{\mathcal P}\|_\square>\epsilon$, there exist subsets $S,T\subset[0,1]$ such that $|\int_{S\times T}(W-W_{\mathcal P})|>\epsilon$. Let $\mathcal P'$ be the refinement of $\mathcal P$ by introducing $S$ and $T$, and that gives at most $4$ sub-parts each.

Define $\langle W,U\rangle$ to be $\int WU$. We know that $\langle W_{\mathcal P},W_{\mathcal P}\rangle=\langle W_{\mathcal P'},W_{\mathcal P}\rangle$ because $W_{\mathcal P}$ is constant on each step of $\mathcal P$, and $\mathcal P'$ is a refinement of $\mathcal P$. Thus, $\langle W_{\mathcal P'}-W_{\mathcal P},W_{\mathcal P}\rangle=0$. By Pythagorean Theorem,
\[\|W_{\mathcal P'}\|_2^2=\|W_{\mathcal P'}-W_{\mathcal P}\|_2^2+\|W_{\mathcal P}\|_2^2>\|W_{\mathcal P}\|_2^2+\epsilon^2,\]
where the latter inequality comes by the Cauchy--Schwarz inequality,
\[\|1_{S\times T}\|_2\|W_{\mathcal P'}-W_{\mathcal P}\|_2\geq|\langle W_{\mathcal P'}-W_{\mathcal P},1_{S\times T}\rangle|=|\langle W-W_{\mathcal P},1_{S\times T}\rangle|>\epsilon.\]
\end{proof}
\begin{proposition}[5.35]
For any $\epsilon>0$, graphon $W$, and $\mathcal P_0$ partition of $[0,1]$, there exists partition $\mathcal P$ refining part of $\mathcal P_0$ into no more than $4^{1/\epsilon^2}$ parts, such that $\|W-W_{\mathcal P}\|_\square\leq\epsilon$.
\end{proposition}
This proposition specifically tells us that starting with any given partition, the regularity argument still works.
\begin{proof}
We repeatedly apply Lemma 5.34 to obtain $\mathcal P_0,\mathcal P_1,\ldots$ partitions of $[0,1]$. For each step, we either have $\|W-W_{\mathcal P}\|_\square\leq\epsilon$ and thus stop, or we know $\|W_{\mathcal P'}\|_2^2>\|W_{\mathcal P}\|_2^2+\epsilon^2$. Because $\|W_{\mathcal P_i}\|_2^2\leq1$, we are guaranteed to stop after fewer than than $\epsilon^{-2}$ steps. We also know that each part is subdivided into no more than $4$ parts at each step, obtaining $4^{\epsilon^{-2}}$ as we desire.
\end{proof}
\begin{theorem}[5.41, Compactness of the space of graphons]
The metric space $(\widetilde{\mathcal W}_0,\delta_\square)$ is compact.
\end{theorem}
\begin{proof}
As $\widetilde{\mathcal W}_0$ is a metric space, it suffices to prove sequential compactness. Fix a sequence $W_1,W_2,\ldots$ of graphons. We want to show that there is a subsequence which converges (with respect to $\delta_\square$) to some limit graphon.

For each $n$, apply the weak regularity lemma (Theorem 5.31) repeatedly, to obtain a sequence of partitions $\mathcal P_{n,1},\mathcal P_{n,2},\mathcal P_{n,3},\ldots$ such that
\begin{enumerate}[label=(\alph*)]
\item $\mathcal P_{n,k+1}$ refines $\mathcal P_{n,k}$ for all $n,k$,
\item $|\mathcal P_{n,k}|=m_k$ where $m_k$ is a function of only $k$, and
\item $\|W_n-W_{n,k}\|_\square\leq1/k$ where $W_{n,k}=(W_n)_{\mathcal P_{n,k}}$.
\end{enumerate}
The weak regularity lemma only guarantees that $|\mathcal P_{n,k}|\leq m_k$, but if we allow empty parts then we can achieve equality.

Initially, each partition may be an arbitrary measurable set. However, for each $n$, we can apply a measure-preserving bijection $\varphi$ to $W_{n,1}$ and $\mathcal P_{n,1}$ so that $\mathcal P_{n,1}$ is a partition of $[0,1]$ into intervals. For each $k\geq2$, assuming that $\mathcal P_{n,k-1}$ is a partition of $[0,1]$ into intervals, we can apply a measure-preserving bijection to $W_{n,k}$ and $\mathcal P_{n,k}$ so that $\mathcal P_{n,k}$ is a partition of $[0,1]$ into intervals, and refines $\mathcal P_{n,k-1}$. By induction, we therefore have that $\mathcal P_{n,k}$ consists of intervals for all $n,k$. Properties (a) and (b) above still hold. While property (c) may not hold, and it's no longer true that $W_{n,k}=(W_n)_{\mathcal P_{n,k}}$, we still know that $\delta_\square(W_n,W_{n,k})\leq1/k$ for all $n,k$. This will suffice for our purposes.

Now, the crux of the proof is a diagonalization argument in countably many steps. Starting with the sequence $W_1,W_2,\ldots$, we will repeatedly pass to a subsequence. In step $k$, we pick a subsequence $W_{n_1},W_{n_2},\ldots$ such that:
\begin{enumerate}
\item the endpoints of the parts of $\mathcal P_{n_i,k}$ all individually converge as $i\to\infty$, and
\item $W_{n_i,k}$ converges pointwise almost everywhere to some graphon $U_k$ as $i\to\infty$.
\end{enumerate}
There is a subsequence satisfying (1) since each partition $\mathcal P_{n,k}$ has exactly $m_k$ parts, and each part has length in $[0,1]$. So consider a subsequence $(W_{a_i})_{i=1}^\infty$ satisfying (1). Each $W_{a_i,k}$ can be naturally identified with a function $f_{a_i,k}:[m_k]^2\to[0,1]$. The space of such functions is bounded, so there is a subsequence $(f_{n_i})_{i=1}^\infty$ of $(f_{a_i})_{i=1}^\infty$ converging to some $f:[m_k]^2\to[0,1]$. Now $f$ corresponds to a graphon $U_k$ which is the limit of the subsequence $(W_{n_i})_{i=1}^\infty$. Thus, (2) is satisfied as well.

To conclude step $k$, the subsequence is relabeled as $W_1,W_2,\ldots$ and the discarded terms of the sequence are ignored. The corresponding partitions are also relabeled. Without loss of generality, in step $k$ we pass to a subsequence which contains $W_1,\ldots,W_k$. Thus, the end result of steps $k=1,2,\ldots$ is an infinite sequence with the property that $(W_{n,k})_{n=1}^\infty$ converges pointwise almost everywhere (a.e.) to $U_k$ for all $k$:
\[
\begin{array}{c|ccccc}
& W_1&W_2&W_3&\cdots&\\\hline
k=1&W_{1,1}&W_{2,1}&W_{3,1}&\cdots&\to U_1\text{ pointwise a.e.}\\
k=2&W_{1,2}&W_{2,2}&W_{3,2}&\cdots&\to U_2\text{ pointwise a.e.}\\
k=3&W_{1,3}&W_{2,3}&W_{3,3}&\cdots&\to U_3\text{ pointwise a.e.}\\
\vdots&\vdots&\vdots&\vdots&&
\end{array}
\]
Similarly, $(\mathcal P_{n,k})_{n=1}^\infty$ converges to an interval partition $\mathcal P_k$ for all $k$.

By property (a), each partition $\mathcal P_{n,k+1}$ refines $\mathcal P_{n,k}$, which implies that $W_{n,k}=(W_{n,k+1})_{\mathcal P_{n,k}}$. Taking $n\to\infty$, it follows that $U_k=(U_{k+1})_{\mathcal P_k}$. Now each $U_k$ can be thought of as a random variable on probability space $[0,1]^2$. From this view, the equalities $U_k=(U_{k+1})_{\mathcal P_k}$ exactly imply that the sequence $U_1,U_2,\ldots$ is a martingale.

The range of each $U_k$ is contained in $[0,1]$, so the martingale is bounded. By the martingale convergence theorem (Theorem 5.39), there exists a graphon $U$ such that $U_k\to U$ pointwise almost everywhere as $k\to\infty$.

Recall that our goal was to find a convergent subsequence of $W_1,W_2,\ldots$ under $\delta_\square$. We have passed to a subsequence by the above diagonalization argument, and we claim that it converges to $U$ under $\delta_\square$. That is, we want to show that $\delta_\square(W_n,U)\to0$ as $n\to\infty$. This follows from a standard ``3-epsilons argument'': let $\epsilon>0$. Then there exists some $k>3/\epsilon$ such that $\|U-U_k\|_1<\epsilon/3$, by pointwise convergence and the dominated convergence theorem. Since $W_{n,k}\to U_k$ pointwise almost everywhere (and by another application of the dominated convergence theorem), there exists some $n_0\in\N$ such that $\|U_k-W_{n,k}\|_1<\epsilon/3$ for all $n>n_0$. Finally, since we chose $k>3/\epsilon$, we already know that $\delta_\square(W_n,W_{n,k})<\epsilon/3$ for all $n$. We conclude that
\begin{align*}
\delta_\square(U,W_n)&\leq\delta_\square(U,U_k)+\delta_\square(U_k,W_{n,k})+\delta_\square(W_{n,k},W_n)\\
&\leq\|U-U_k\|_1+\|U_k-W_{n,k}\|_1+\delta_\square(W_{n,k},W_n)\leq\epsilon.
\end{align*}
The second inequality uses the general bound that $\delta_\square(W_1,W_2)\leq\|W_1-W_2\|_\square\leq\|W_1-W_2\|_1$ for graphons $W_1,W_2$.
\end{proof}
